% Conclusions: a focused two to three paragraph close that restates the
% contribution and the forward path without introducing new data, referencing the
% Results figures and tables rather than repeating numbers.

An autonomous Claude Code Opus 4.7 (1M context) Max run executed every entry
point of the generated platform, passed the full automated test suite, drove the
ten-gate accept, block, and escalate surface summarised in
\autoref{tab:decisionsurface}, cleared every iteration of the deterministic
sweep, processed the featured tournament and the thousand-row sensor stream, and
reproduced all three large artifacts from a single seed, all under one pull
request. That record operationalises the thesis. For a single autonomous surgical
humanoid the assurance process, not code generation, is the substantial and
decision-bearing work: the control behaviors are produced in microseconds, while
turning any one of them into a shippable result demands an exact verification
fraction, an agreement bar that reaches 1.00 on the catastrophe gates, the tight
dispersion bounds shown in \autoref{fig:forest} and \autoref{tab:thresholds}, a
recorded human review, and an extra hard predicate where the risk is highest. The
assurance layer is larger, stricter, and more thoroughly exercised than the code
it judges, which is exactly what the embodiment of a two-handed standing surgeon
requires, and an inexpensive autonomous agent carried that heavier half cheaply
and reproducibly.

Binding every gate to published external standards already used in real life is
what turns that inexpensive run into defensible evidence, because a reviewer can
follow any decision to its governing designation rather than to an invented
threshold, as the binding in \autoref{fig:matrix} makes plain. This checks real
assurance steps off, determinism, traceable standards, pass-or-fail verification,
validation against independent references, recorded human review, and a
hand-back-to-human default, so the field does not have to re-derive them for each
deliverable, and it builds directly on the prior two-stage VVUQ framework that
placed verification over generation \cite{kawchak2026vvuq01} and on the
credibility standards that frame the bounds \cite{asmevv40, nasastd7009a}.
Adapting the approach across oncology clinical trials therefore promises faster,
less expensive, and more patient-beneficial deliverables than conventional
verification.

The opportunity is real and the posture is careful. The kind of fast,
standards-anchored, autonomous assurance demonstrated here is what the move toward
real-time clinical trials invites, and United States leadership in
surgical-humanoid VVUQ tied to external standards would serve patient safety,
efficacy, and trial speed together. At the same time, the Unitree H2-Surgical 1.0
remains a clearly labelled hypothetical 2030 platform whose every candidate
behavior must clear the ten gates and a recorded human reviewer, and whose
deployment would require the full regulatory pathway, before any non-simulated
use. The contribution is the assurance scaffold and the evidence that an
autonomous agent can run it well, which is the part the future of patient safety
most needs settled.
